%==============================================================================
% CHAPTER 5
%==============================================================================
\chapter{Difficulties, false starts and schedule}
\markboth{Chapter 5 — Difficulties, false starts and schedule}{}

This chapter reports the obstacles encountered. It is not a catalogue of
anecdotes: each of these difficulties changed a design decision or a working
method.

\section{Technical difficulties}

\subsection{Hardware acceleration of the language model}

Development was initially slowed by agent invocations exceeding four minutes,
the language model running on the CPU. As the working machine had a capable
graphics card, exploiting it seemed straightforward. This was the
longest-standing difficulty to resolve, for three linked reasons.

First, the usual acceleration software stack for that graphics card vendor does
not cover the model available on the machine under the Linux environment used.
It had to be identified that an alternative route, based on a standard graphics
interface, was available and should be preferred.

Second, a misleading error message cost time: the log stated that the graphics
driver was too old, whereas it was up to date. That message in fact came from a
check relating to the stack that had not been selected, unrelated to the route
actually in use. The lesson is a classic one but worth retaining: an explicit
error message is not necessarily a \emph{relevant} error message.

Third, a second installation of the same service, made by mistake inside the
Linux environment instead of the host system, was silently intercepting every
request. The measurements obtained were therefore those of the slow
configuration, even though the fast one was in place. This confusion illustrates
a methodological point: when two equivalent services share a port, no
measurement is reliable until one has checked \emph{which} of them answers.

\subsection{Compiling the solver}

Installing OFFBEAT ran into a simple obstacle: the repository address given in
the project's internal documentation was no longer valid. The actual public
repository had to be identified, and its build structure turned out to differ
from the one described. Compilation itself then proceeded without error.

This episode led to a more general remark: a project's documentation ages, and
wrong documentation costs more than absent documentation, because it actively
points in the wrong direction.

\subsection{Undeclared dependencies}

Porting to a fresh machine revealed three Python libraries used by the code but
absent from the dependency list. They were present by accident on the original
machine. This kind of defect only shows up during a transfer, which argues for
performing such a transfer \emph{early} rather than at the end of a project.

\section{False starts and reorientations}

\subsection{A correct but ineffective diagnosis}

The most instructive false start concerns the very design of self-healing. The
initial assumption, inherited from the reference project, was that a crash is
sufficiently characterised by its signature in the run log. The stress-testing
of chapter~4 showed that this assumption is insufficient: two crashes with
identical signatures may call for opposite remedies, or for no automatable
remedy at all.

The resulting reorientation is a shift of effort: rather than endlessly
enriching the base of crash patterns, it is more effective to \textbf{validate
inputs before execution}. A case that cannot be physically valid should never
reach the solver.

\subsection{A conservative approximation that was not}

The second false start is the initial choice to measure cladding quantities over
the whole domain, documenting it as a conservative approximation. The reasoning
seemed sound: overestimating a stress can only err on the side of safety.

It was wrong. In the case of hoop stress, the extremum retained lay inside the
pellet and carried a sign opposite to the actual stress in the cladding —
tension reported where the cladding was in compression. An approximation
prudent in principle had turned into a qualitatively false result.

The lesson drawn goes beyond this particular case: a choice documented as
conservative must be \emph{verified} to be so, not assumed. The correction was
validated on a real case, where it did change the strain result by a factor of
five and reversed the sign of the stress.

\subsection{A rigorous validation of an uninteresting problem}

The third false start is the one from which I learned the most, because it
concerned not a bug but a \emph{method}.

The surrogate had been validated with what seemed to me every precaution:
\textit{leave-one-out} cross-validation, chosen precisely because it is more
honest than a goodness-of-fit coefficient, and coefficients of determination
above $0.998$ on the three quantities tracked. The statistical method was
correct and the result excellent.

The problem is that the question asked was not. The sampling domain — a few
thousand seconds — had been chosen so that the design of experiments would
remain quick to build, at a time when the solver was not compiled locally. Yet
over that domain the rod reaches thermal equilibrium: the irradiation duration,
the model's second variable, carried almost no information. I had therefore
rigorously validated a model on an almost one-dimensional problem, while the
dashboard was presenting its predictions as two-parameter safety margins.

What made this visible was not a metric but a confrontation: the reference
simulation ran to \SI{6.3e7}{\second} when the surrogate had never seen beyond
\SI{6500}{\second}. A ratio of ten thousand between the training domain and the
domain of use cannot be read off an $R^2$.

The lesson is twofold. On the one hand, a validation metric only qualifies the
fit to a dataset, never the relevance of that dataset: a model can be
impeccably validated and nonetheless useless. On the other hand, the practical
constraint that had dictated the choice — the cost of building the design of
experiments — had disappeared once the solver was compiled locally, without my
re-examining the decision it had motivated. Design assumptions deserve to be
re-evaluated when the constraints that justified them change.

\subsection{On the use of a programming assistant}

Development was carried out with the support of a programming assistant based on
a language model. The perspective gained is worth recording, since the subject
of the internship concerns this very class of tools.

The assistance proved very effective for producing structured code and for
exploring an unfamiliar codebase. It proved, on the other hand, incapable of
detecting on its own the correctness defects of table~\ref{tab:bugs}: these were
found only by \emph{confronting numerical results with a physical expectation} —
why is the cladding temperature equal to the pellet centreline temperature? why
is a closed gap reported as safe? Physical judgement remains the decisive link,
and this observation applies equally to the agent developed here.

\clearpage

%==============================================================================
% SCHEDULE (full page just before the conclusion, per the guidelines)
%==============================================================================
\section*{Internship schedule}
\addcontentsline{toc}{section}{Internship schedule}

\vspace{0.4cm}

\noindent
\textit{\color{red}[To be adjusted to the actual internship dates. The phases
below reflect the order in which the work was actually built.]}

\vspace{0.6cm}

\begin{table}[htbp]
\centering
\caption[Internship schedule by phase]{Progress of the internship by phase.}
\label{tab:planning}
\small
\begin{tabularx}{\textwidth}{@{}>{\raggedright\arraybackslash}p{2.6cm}
                                >{\raggedright\arraybackslash}p{4.3cm}X@{}}
\toprule
\textbf{Period} & \textbf{Phase} & \textbf{Deliverable / milestone}\\
\midrule
Weeks 1--2 & Familiarisation
  & Fuel physics, first steps with OpenFOAM and OFFBEAT, study of the
    reference project AutoFLUKA\\
\addlinespace
Weeks 3--4 & Software foundation
  & Project tree, provider-agnostic language model factory, development
    environment\\
\addlinespace
Weeks 5--6 & Execution and self-healing
  & Solver execution tool validated on a real case; error pattern base and
    repair loop\\
\addlinespace
Weeks 7--8 & Complete agent
  & Supervisor, case generation tool, post-processing; first end-to-end
    chaining\\
\addlinespace
Weeks 9--10 & Interface and documentation
  & Web interface, documentary assistant based on semantic search\\
\addlinespace
Weeks 11--13 & Digital twin
  & Safety analyser, prognosis, assimilation, dashboard, Gaussian-process
    surrogate\\
\addlinespace
Weeks 14--15 & Validation and consolidation
  & Reference simulation over a complete cycle, stress-testing of self-healing,
    correction of the correctness defects\\
\addlinespace
Week 16 & Writing up
  & Report and defence\\
\bottomrule
\end{tabularx}
\end{table}

\clearpage
